\documentclass[a4paper, 12pt]{report}

%====================== PACKAGES ======================

\usepackage[french]{babel}
\usepackage[utf8x]{inputenc}
%pour gérer les positionnement d'images
\usepackage{float}
\usepackage{amsmath}
\usepackage{graphicx}
\usepackage[colorinlistoftodos]{todonotes}
\usepackage{url}
%pour les informations sur un document compilé en PDF et les liens externes / internes
\usepackage{hyperref}
%pour la mise en page des tableaux
\usepackage{array}
\usepackage{tabularx}
%pour utiliser \floatbarrier
%\usepackage{placeins}
%\usepackage{floatrow}
%espacement entre les lignes
\usepackage{setspace}
%modifier la mise en page de l'abstract
\usepackage{abstract}
%police et mise en page (marges) du document
\usepackage[T1]{fontenc}
\usepackage[top=1.5cm, bottom=2cm, left=1.5cm, right=1.5cm]{geometry}
%Pour les galerie d'images
\usepackage{subfig}
\usepackage{eurosym}

%====================== INFORMATION ET REGLES ======================

%rajouter les numérotation pour les \paragraphe et \subparagraphe
\setcounter{secnumdepth}{4}
\setcounter{tocdepth}{4}

\hypersetup{							% Information sur le document
pdfauthor = {Sugdenaz EKICI,
			Yahia KHERZA,
			Olivier MARAVAL,
    		Valentin VIRET-JACQUOT},			% Auteurs
pdftitle = {1_Etude},			% Titre du document
pdfsubject = {Mémoire de Projet},		% Sujet
pdfkeywords = {Tag1, Tag2, Tag3, ...},	% Mots-clefs
pdfstartview={FitH}}					% ajuste la page à la largueur de l'écran
%pdfcreator = {MikTeX},% Logiciel qui a crée le document
%pdfproducer = {}} % Société avec produit le logiciel

%======================== DEBUT DU DOCUMENT ========================

\begin{document}

\addtocontents{toc}{\protect\thispagestyle{empty}}

%régler l'espacement entre les lignes
\newcommand{\HRule}{\rule{\linewidth}{0.5mm}}

%page de garde
\begin{titlepage}
\begin{center}

% Upper part of the page. The '~' is needed because only works if a paragraph has started.
\includegraphics[width=0.5\textwidth]{./images/InfoLogoQuadriH.png}~\\[1cm]

\textsc{\LARGE SAE 1.03 - BUT INFORMATIQUE - GROUPE 1}\\[1.5cm]

\textsc{\Large }\\[0.5cm]

% Title
\HRule \\[0.4cm]

{\huge \bfseries Architecture logicielle et notice d'installation / utilisation du poste de développement\\[0.4cm] }

\HRule \\[1.5cm]

% Author and supervisor
\begin{minipage}{0.4\textwidth}
\begin{flushleft} \large
\emph{Auteur:}\\
Sugdenaz \textsc{Ekici}(\textit{A1})\\
Yahia \textsc{Kherza}(\textit{A1})\\
Olivier \textsc{Maraval}(\textit{A1})\\
Valentin \textsc{Viret-Jacquot}(\textit{A1})
\end{flushleft}
\end{minipage}
\begin{minipage}{0.4\textwidth}
\begin{flushright} \large
\emph{Client:} \\
Michel \textsc{Salomon}\\
\emph{Référent:} \\
Olivier \textsc{Maraval}
\end{flushright}
\end{minipage}

\vfill

% Bottom of the page
{\large \today}

\end{center}
\end{titlepage}

%page blanche
\newpage
~
\thispagestyle{empty}


\pagenumbering{gobble}
\tableofcontents
\thispagestyle{empty}

%ne pas numéroter le sommaire

\newpage

%espacement entre les lignes d'un tableau
\renewcommand{\arraystretch}{1.5}

%====================== INCLUSION DES PARTIES ======================

~
\thispagestyle{empty}


\newpage
\pagenumbering{arabic}

\setcounter{page}{5}

\chapter{Votre demande}

\section{Résumé de votre demande}

Votre requête concerne une configuration de travail personnalisée pour le développement d'applications web. 

Vous nous avez spécifié un certain nombre d'outils et d'environnements de développement que vous souhaitez avoir à disposition, notamment : 
\begin{itemize}
\item MariaDB / MySQL pour les bases de données, 
\item Python3 avec la possibilité d'avoir des environnements virtuels, Flask et un serveur web.
\end{itemize}

\section{Notre méthode}

Au vue de votre demande plutôt ouverte, notre équipe vous a détaillé différentes possibilités dans la suite de ce document et vous a guidé vers l’option qui lui semble la plus adaptée.

\subsection{Concernant l'architecture matérielle}

Nous vous avons proposé quatre architectures matérielles possibles : deux pour un poste de travail portable et deux pour un poste de travail fixe. Vous avez à chaque fois une option économique et une option confort.

Nous avons inclus dans nos propositions les périphériques d'usage habituels tels qu'un clavier, une souris, un écran, etc. De même pour les ordinateurs portables afin de pouvoir les utiliser comme poste fixe.

Le prix de chaque composant est indiqué. Nous avons intégré un lien web vers la description de chaque composant à toute fins utiles. Vous trouverez également un coût prévisionnel pour chaque configuration.

\subsection{Concernant l'architecture logicielle}

\subsubsection{Le système d'exploitation}

Pour des raisons de sécurité, de flexibilité et de compatibilité, notre équipe s’est arrêtée sur le système d’exploitation Linux. 

Nous vous avons proposé différentes distributions et environnements de bureau afin que vous puissiez rapidement prendre en main votre nouveau poste de travail.

\subsubsection{La suite logicielle}

Notre équipe a décidé de vous proposer une suite logicielle complète vous permettant de commencer immédiatement le développement d’applications web.

Chacun des logiciels que nous avons sélectionné s’accompagne d’une courte description vous indiquant son utilité.

Pour chaque logiciel payant de cette liste nous vous avons indiqué une alternative gratuite en vous indiquant les inconvénients potentiels à un tel choix. Dans le cas où les fonctionnalités perdues s'avéraient anecdotiques, nous avons uniquement inclus l’alternative gratuite.

\chapter{Propositions d'architecture matérielle}

\section{Ordinateur Fixe}

\subsection{Budget 1}

\begin{itemize}
	\item Processeur https://www.ldlc.com/fiche/PB00497548.html 104€95
	\item Carte mère https://www.ldlc.com/fiche/PB00346172.html 119€95
	\item Ram https://www.ldlc.com/fiche/PB00191551.html 58€94
	\item Carte Graphique https://www.ldlc.com/fiche/PB00462561.html 79€96
	\item SSD https://www.ldlc.com/fiche/PB00402115.html 44€95
	\item HDD https://www.ldlc.com/fiche/PB00254175.html 72€95
	\item Boitier https://www.ldlc.com/fiche/PB00275652.html 99.95
	\item Alimentation https://www.ldlc.com/fiche/PB00548770.html 164€95
	\item Graveur DVD https://www.ldlc.com/fiche/PB00213964.html 29€95
	\item Carte réseau https://www.ldlc.com/fiche/PB00487856.html 39€95
	\item Câble SATA https://www.ldlc.com/fiche/PB00016236.html
	\item Ventilateur https://www.ldlc.com/fiche/PB00264734.html
	\item Souris https://www.ldlc.com/fiche/PB00243882.html
	\item Clavier https://www.ldlc.com/fiche/PB00242405.html
	\item Ecran https://www.ldlc.com/fiche/PB00233120.html
	\item Haut parleur https://www.ldlc.com/fiche/PB00154969.html
	\item Webcam https://www.ldlc.com/fiche/PB00213751.html
	\item Clé Usb https://www.ldlc.com/fiche/PB00547743.html
	\item Tapis de Souris https://www.ldlc.com/fiche/PB00466685.html
	\item Casque https://www.ldlc.com/fiche/PB00366138.html
	\item PRIX 1265,06\euro
\end{itemize}

\subsection{Budget 2}

\begin{itemize}
	\item Processeur https://www.ldlc.com/fiche/PB00497963.html
	\item Carte mère https://www.ldlc.com/fiche/PB00345955.html
	\item Mémoire https://www.ldlc.com/fiche/PB00191544.html
	\item Carte graphique https://www.ldlc.com/fiche/PB00543736.html 
	\item SSD https://www.ldlc.com/fiche/PB00265389.html
	\item HDD https://www.ldlc.com/fiche/PB00254175.html
	\item Boitier https://www.ldlc.com/fiche/PB00278344.html
	\item Alimentation https://www.ldlc.com/fiche/PB00401724.html
	\item Carte réseau (Usb) https://www.ldlc.com/fiche/PB00266161.html
	\item Cable SATA https://www.ldlc.com/fiche/PB00016236.html
	\item Ventilateur https://www.ldlc.com/fiche/PB00264734.html
	\item Souris https://www.ldlc.com/fiche/PB00243882.html
	\item Clavier https://www.ldlc.com/fiche/PB00103723.html
	\item Ecran https://www.ldlc.com/fiche/PB00543634.html
	\item Enceinte https://www.ldlc.com/fiche/PB00101582.html
	\item Webcam https://www.ldlc.com/fiche/PB00213751.html
	\item Clé Usb https://www.ldlc.com/fiche/PB00547743.html
	\item Tapis de Souris https://www.ldlc.com/fiche/PB00466685.html
	\item Casque https://www.ldlc.com/fiche/PB00351651.html
	\item PRIX 915,65\euro 
\end{itemize}

\subsection{Budget 3}

\section{Ordinateur Portable}

\subsection{Budget 1}


\subsection{Budget 2}
\chapter{Proposition d'architecture logicielle}

\section{Système d'exploitation}

\begin{table}[ht]
	\centering
	\newcolumntype{C}{>{\centering\arraybackslash}p{3cm}}
	\begin{tabularx}{\textwidth}{|C|X|X|}
		\hline
		\textbf{Distribution} & \textbf{Avantages} & \textbf{Inconvénient} \\
		\hline
		\href{https://ubuntu.com}{\textbf{Ubuntu}\vspace{5px} \centering \includegraphics[width=0.7\linewidth]{images/Ubuntu}} & Dernières mises à jour des logiciels ; inclus dans ses dépots une sélection de logiciel et drivers propriétaire & Peut s'avérer moins stable mais suffisamment fiable pour une station de travail. \\
		\hline
		\href{https://www.debian.org}{\textbf{Debian}\vspace{5px} \centering \includegraphics[width=0.7\linewidth]{images/Debian}} & Globalement plus stable ce qui peut-être intéressant pour un serveur & Installer un logiciel propriétaire peut s'avérer plus compliqué ; les drivers open-sources peuvent impacter les performances de certains composants \\
		\hline
	\end{tabularx}
	\caption{Avantages et inconvénients des distributions}
	\label{tab:distrib}
\end{table}


\begin{table}[ht]
	\centering
	\newcolumntype{C}{>{\centering\arraybackslash}p{4cm}}
	\begin{tabularx}{\textwidth}{|C|X|X|}
		\hline
		\textbf{Environnement de Bureau} & \textbf{Description} \\
		\hline
		\href{https://mate-desktop.org/fr/}{\textbf{Mate} \linebreak \vspace{5px} \centering \includegraphics[height=0.5\linewidth]{images/mate}} & Un environnement de bureau léger et modulable qui vous permettra d'exploiter au maximum les performances de votre machine. Son interface peut se réveler austère (voir figure) \\
		\hline
		\href{https://kde.org/fr/}{\textbf{KDE Plasma 5} \linebreak \vspace{5px} \centering \includegraphics[width=0.5\linewidth]{images/kdelogo}} & Un environnement de bureau complet, personnalisable et agencé comme Windows. Il est légèrement plus lourd et encombré que Mate mais sera beaucoup plus confortable pour un usage quotidien. (voir figure) \\
		\hline
	\end{tabularx}
	\caption{Présentation des environnements de bureau}
	\label{tab:desktop}
\end{table}

\section{Organisation du système}

\subsection{Les comptes utilisateurs}

Un compte administrateur \textbf{"root"} sera disponible. Il ne doit être utilisé que par une personne compétente dans le cas d'un problème système à régler.\\

Nous vous créerons un compte utilisateur qui vous permettra de travailler en sécurité sur votre machine. Ce compte aura accès à la commande \textit{\textbf{sudo}} pour vous permettre l'installation de nouvelles applications ou la mise à jour de votre système.

Un second compte \textbf{"invité"} sera crée pour une personne souhaitant utiliser votre station de travail à des fins bureautique par exemple. Il n'aura pas accès à la commande \textit{\textbf{sudo}} pour des raisons de sécurité.

Vous aurez la possibilité d'ajouter un compte utilisateur pour chaque personne souhaitant utiliser votre station de travail (nous vous recommandons de ne pas lui donner l'accès à la commande \textit{\textbf{sudo}} qui lui accorderait les droits d'administration sur votre système).

\subsection{Organisation du disque}

\subsubsection{Partitionnement}

Nous n'effectueront pas de double-boot dans la mesure où Windows ne vous sera pas utile pour le développement web, nous nous contenterons donc d'une simple installation de Linux pour des raisons d'optimisation de stockage. Dans le cas où vous choisiriez un ordinateur tournant sur MacOS, nous laisserions le système d'exploitation d'origine pour des raisons de compatibilité et de confort utilisateur.\\

Dans le cas de l'installation d'un Linux, nous partitionnerons le disque de la façon suivante :
\begin{itemize}
	\item \textit{/boot} : 300 Mio - Cette partition contiendra tout ce qui permet à l'OS de booter.
	\item \textit{/} : Occupera l'espace restant du disque. Elle contiendra tous les fichiers de l'OS.
\end{itemize}
\vspace{10px}
Nous n'avons pas inclus de partition \textbf{SWAP} puisque toutes nos configurations possèdent au moins 16Go de mémoire vive, vous ne devriez donc avoir aucun problème de stabilité. Par ailleurs, la fonction d'hibernation nécessiterait un swap au minimum équivalent à votre mémoire vive ce qui viendrait réduire l'espace de stockage dont vous disposez.

\subsection{Arborescence des Fichiers}
Chaque utilisateur aura son dossier dans \textit{/home/<nom-de-l'utilisateur>} dans lequel se trouvera son travail. Ce dossier ne sera pas accessible par les autres utilisateurs du système. Il trouvera à l'intérieur :
\begin{itemize}
	\item Un dossier \textbf{"projets"} ou placer ses travaux en cours, c'est dans ce dossier qu'il pourra cloner les dépôts git sur lesquels il collaborera.
	\item Un dossier \textbf{"archives"} ou poser ses travaux finis, pour une éventuelle réutilisation par exemple.
	\item Un dossier \textbf{"ressources"} dans lequel il pourra placer tous les fichiers qu'il utilisera régulièrement dans son travail.
\end{itemize}

Un dossier \textit{/home/public} sera accessible en lecture et en écriture par tous les utilisateurs du système pour le partage de certains fichiers. Aucun utilisateur ne disposera du droit d'execution dans ce dossier pour des raisons de sécurité.

\section{Logiciels}

\subsection{Logiciels essentiels}\label{3.3.1}

\begin{itemize}
	\item \textbf{Mozilla Firefox, Google Chrome et GNOME Web} : pour pouvoir tester vos sites sur un navigateur Gecko, Blink et Webkit.
	\item \textbf{Visual Studio Code} : un éditeur de code polyvalent permettant la programmation dans de multiple langages.
	\item \textbf{DBeaver} : un IDE SQL open-source vous permettant de gérer plus facilement votre base de données.
	\item \textbf{Python3 avec venv et Flask} : pour vous permettre de créer des applications web avec la possibilité d'avoir un environnement virtuel.
	\item \textbf{Apache} : le serveur web le plus utilisé au monde 
	\item \textbf{MariaDB} : comme système de gestion de base de données
	\item \textbf{Zoom, Discord et/ou Slack} pour la communication entre clients et collaborateurs.
\end{itemize}

\subsection{Options supplémentaires gratuites}\label{3.3.2}

\begin{itemize}
	\item \textbf{Pycharm} : un IDE python efficace pour le développement d'application web sous flask, plus complet que Visual Studio Code mais moins polyvalent.
	\item \textbf{Looping} : Un outil de création de MCD essentiel pour vous assister dans la conception de votre base de données.
	\item \textbf{git} : pour une meilleure gestion de votre projet et une éventuelle collaboration sur ce dernier.
	\item \textbf{GIMP} : pour la retouche d'images.
	\item \textbf{LibreOffice} : une suite bureautique intéressante si vous prévoyez de sortir des devis, factures, suivre vos dépenses, etc.
\end{itemize}


\subsection{Alternatives Payantes}

\begin{itemize}
	\item \textbf{DataGrip} (en remplacement de DBeaver pour 229\euro\ HT par an) : Plus complet et ergonomique que DBeaver c'est une option intéressante si vous comptez travailler sur des bases de données assez complexes.
	\item \textbf{WebStorm} (pour 159\euro\ HT par an) : Un IDE spécialisé dans le javascript. Intéressant si vous souhaitez concevoir des sites plus créatifs faisant usage intensif du javascript.
\end{itemize}

\chapter{Préconisation d'un poste de développement}

\section{Comparaison des propositions précédentes}

\subsection{Type de poste de travail}

Un ordinateur fixe offre généralement une meilleure performance et une plus grande capacité de stockage que les ordinateurs portables pour un même prix. En revanche, si vous prévoyez des déplacements réguliers, il peut être intéressant de vous tourner vers un ordinateur portable quitte à sacrifier un peu de puissance ou augmenter votre budget.

Dans votre cas, une grande partie de votre travail ne nécessite pas de vous déplacer chez vos client, aussi, nous vous conseillons de vous tourner vers un ordinateur fixe.\\

\textbf{Notre conseil :} Un poste de travail fixe

\subsection{Performances recommandées}

Parmi les propositions que nous vous avons faites, et d'après votre situation actuelle, nous vous conseillons d'éviter l'option la plus économique. En effet, les économies que vous effectuerez à l'achat se répercuteront sur votre productivité.

En tant que développeur web, vous serez amené à effectuer de la retouche d'image voir de la conversion vidéo à des fins d'intégration sur vos pages web. L'absence de carte graphique rendra cette partie du travail lente et pénible et rallongera le temps de développement de vos sites ou en diminuera leur qualité.

Si vous deviez choisir l'option économique que nous vous avons proposé, elle s'avèrera tout de même assez puissante pour vous permettre de travailler confortablement sur du développement web ne demandant pas de calcul graphiques poussés.\\

\textbf{Notre conseil :}  Le poste de travail fixe intermédiaire \ref{2.1.2}

\subsection{Choix du système d'exploitation et de l'environnement de bureau}

Au vue du peu de différence notable entre Debian et Ubuntu, on choisira la distribution la plus simple d'accès pour vous à savoir Ubuntu.

De la même façon, on choisira d'installer l'environnement KDE Plasma 5 pour sa facilité d'utilisation et sa modularité. Il inclue aussi une suite logicielle intéressante ainsi qu'un trousseau de clé qui vous permettra d'activer, notamment, les fonctions de synchronisation de Visual Studio Code sans plus de configuration.

Il existe une déclinaison d'Ubuntu nommé Kubuntu qui inclue cet environnement de bureau, c'est donc celle-ci que nous allons installer.

\textbf{Notre conseil :} Kubuntu - Kde Plasma 5

\subsection{Logiciels conseillés}

Parmi les logiciels que nous vous avons proposé, nous vous conseillons d'installer l'intégralité des logiciels gratuits. En effet, les alternatives payantes n'ont d'intérêt que si vous savez déjà pourquoi vous souhaiteriez les utiliser.

\textbf{Notre conseil :} L'intégralité des logiciels détaillés en \ref{3.3.1} et \ref{3.3.2}

\subsection{Solutions de stockage}

Bien qu'un simple disque dur soit une alternative économique intéressante, nous pensons que la liberté et l'espace de stockage offerte par un NAS justifie son inclusion dans votre configuration. En effet, il sera plus simple pour vous de proposer un espace de stockage accessible en permanence et ceux peut importe que vous soyez à votre bureau ou en déplacement.

Par ailleurs, nous vous recommandons vivement un système de stockage de type cloud pour sécuriser vos fichiers les plus importants en cas de panne de votre ordinateur ou de votre stockage local.

\textbf{Notre conseil :} Un NAS et deux HDD 4TO + Dropbox

\section{Prix}

\begin{table}[ht]
	\centering
	\newcolumntype{C}{>{\centering\arraybackslash}p{3cm}}
	\begin{tabularx}{\textwidth}{|C|X|}
		\hline
		\textbf{Matériel} & 1351.13\euro \\
		\hline
		\textbf{Logiciel} & 0\euro \\
		\hline
		\textbf{Stockage} & 502.83\euro + 192\euro/an \\
		\hline
		\hline
		\textbf{Total} & 1853.96\euro + 192\euro/an \\
		\hline		
	\end{tabularx}
	\caption{Prix total de la configuration conseillée}
	\label{tab:prix}
\end{table}

\newpage

%récupérer les citation avec "/footnotemark"
\nocite{*}



\end{document}